El presente trabajo expone la evolución del sistema \textit{FireSafe} hacia una arquitectura de Internet de las Cosas (IoT) distribuida, diseñada para el monitoreo ambiental crítico y la detección temprana de riesgos. A diferencia de su predecesor, esta versión implementa una topología de red en la que múltiples nodos periféricos basados en el microcontrolador ESP32 realizan el procesamiento en el borde (\textit{Edge Computing}), gestionando la concurrencia de tareas mediante el sistema operativo de tiempo real FreeRTOS \citep{valvano2014volume1}. 

La infraestructura se apoya en un protocolo de mensajería asíncrona MQTT bajo un modelo de publicación/suscripción, garantizando una comunicación eficiente y escalable entre los nodos y el servidor central en una Raspberry Pi. El sistema integra un \textit{backend} desarrollado en Django para la persistencia de datos y un \textit{dashboard} interactivo en React que permite la visualización en tiempo real y la gestión de alertas. Los resultados demuestran que el análisis de circuitos propuesto por \citet{boylestad2009} y la implementación de interfaces de tiempo real \citep{valvano2014volume2} permiten una respuesta inmediata ante fluctuaciones térmicas o presencia de gases combustibles, validando la viabilidad de soluciones industriales de bajo costo.

\textbf{Palabras clave:} IoT distribuido, MQTT, FreeRTOS, monitoreo ambiental, sistemas de tiempo real.